% 408 历年真题索引（2009-2025）
% 按学科→章节→知识点组织，标注 年份-题号-题型-分值

\chapter{408 历年真题索引（2009--2025）}

\noindent 本附录将 2009--2025 年全部 408 统考真题按本书章节和知识点分类，注明题号、题型和分值。建议在每章学习后，对照本索引完成对应年份真题的训练。

\section{数据结构}

\subsection*{第1章 线性表}

\begin{tablebox}{线性表 — 真题索引}
\centering
\begin{tabular}{|c|l|c|c|c|}
\hline
年份 & 知识点 & 题号 & 题型 & 分值 \\
\hline
2009 & 单链表查找倒数第k个结点 & 42 & 综合（代码）& 15 \\
2010 & 顺序表循环左移 & 42 & 综合（代码）& 15 \\
2011 & 顺序表求中位数（两个等长升序序列）& 42 & 综合（代码）& 15 \\
2012 & 单链表查找公共后缀起始位置 & 42 & 综合（代码）& 15 \\
2013 & 顺序表主元素（出现次数 $>n/2$）& 41 & 综合（代码）& 13 \\
2014 & 顺序表删除重复元素（有序表）& 41 & 综合 & 13 \\
2015 & 单链表删除绝对值相等的结点 & 41 & 综合（代码）& 15 \\
2016 & 顺序表划分（枢轴划分）& 43 & 综合 & 15 \\
2017 & 顺序表三元组最小距离 & 41 & 综合（代码）& 15 \\
2018 & 顺序表未出现的最小正整数 & 41 & 综合（代码）& 13 \\
2019 & 单链表就地逆置（带头结点）& 41 & 综合（代码）& 13 \\
2020 & 三元组最小距离（变体）& 41 & 综合 & 13 \\
2021 & 顺序表逆置综合 & 41 & 综合 & 13 \\
2022 & 单链表删除指定范围元素 & 41 & 综合（代码）& 13 \\
2023 & 顺序表找出两等长有序序列合并后的中位数 & 41 & 综合 & 13 \\
2024 & 链表拆分（将链表按奇偶位置拆分）& 41 & 综合（代码）& 13 \\
2025 & 链表删除重复结点 & 41 & 综合（代码）& 13 \\
\hline
\end{tabular}
\end{tablebox}

\noindent\textbf{高频命题规律：}
\begin{itemize}
    \item \textbf{顺序表}：循环移位、主元素、未出现的最小正整数、中位数——都是 $O(n)$ 时间 $O(1)$ 空间的经典套路。
    \item \textbf{链表}：逆置、删除特定结点、查找公共结点——快慢指针和 dummy 头结点是核心技巧。
    \item 2009-2025 共 17 年，\textbf{每年至少一道线性表综合题（代码题）}，是 DS 中最高频的大题考点。
\end{itemize}

\subsection*{第2章 栈与队列}

\begin{tablebox}{栈与队列 — 真题索引}
\centering
\begin{tabular}{|c|l|c|c|c|}
\hline
年份 & 知识点 & 题号 & 题型 & 分值 \\
\hline
2009 & 栈的出入序列判定 & 2 & 单选 & 2 \\
2010 & 循环队列空/满判定 & 1 & 单选 & 2 \\
2011 & 栈的出入序列 & 2 & 单选 & 2 \\
2012 & 中缀表达式转后缀 & 2 & 单选 & 2 \\
2013 & 栈在递归中的作用 & 2 & 单选 & 2 \\
2014 & 循环队列 & 1 & 单选 & 2 \\
2016 & 栈与队列的应用 & 2 & 单选 & 2 \\
2017 & 栈的序列 & 2 & 单选 & 2 \\
2018 & 共享栈 & 2 & 单选 & 2 \\
2019 & 队列在层次遍历中的应用 & 3 & 单选 & 2 \\
2021 & 循环队列指针计算 & 3 & 单选 & 2 \\
2023 & 栈的后缀表达式求值 & 1 & 单选 & 2 \\
2024 & 循环队列 & 1 & 单选 & 2 \\
\hline
\end{tabular}
\end{tablebox}

\noindent\textbf{高频：}栈的出入序列判定（几乎每年）、循环队列 front/rear 指针计算（408 特色考点）。

\subsection*{第3章 串}

\begin{tablebox}{串 — 真题索引}
\centering
\begin{tabular}{|c|l|c|c|c|}
\hline
年份 & 知识点 & 题号 & 题型 & 分值 \\
\hline
2012 & KMP 算法 next 数组计算 & 6 & 单选 & 2 \\
2014 & KMP 算法 next 数组 & 6 & 单选 & 2 \\
2015 & KMP next 数组计算 & 8 & 单选 & 2 \\
2016 & KMP nextval 数组 & 8 & 单选 & 2 \\
2018 & KMP next 数组 & 8 & 单选 & 2 \\
2019 & KMP next 数组（下标从0开始）& 8 & 单选 & 2 \\
2020 & KMP next 数组 & 8 & 单选 & 2 \\
2021 & KMP 匹配过程比较次数 & 8 & 单选 & 2 \\
2023 & KMP next 数组 & 6 & 单选 & 2 \\
2024 & KMP next 数组 & 5 & 单选 & 2 \\
2025 & KMP 算法 & 6 & 单选 & 2 \\
\hline
\end{tabular}
\end{tablebox}

\noindent\textbf{规律：}408 对串的考查几乎全部集中在 KMP next 数组的手工计算，每年 1 道选择题（2分）。

\subsection*{第4章 树与二叉树}

\begin{tablebox}{树与二叉树 — 真题索引（选择题）}
\centering
\begin{tabular}{|c|l|c|c|c|}
\hline
年份 & 知识点 & 题号 & 题型 & 分值 \\
\hline
2009 & 平衡二叉树定义 & 4 & 单选 & 2 \\
2010 & 二叉树性质（$n_0=n_2+1$）& 3 & 单选 & 2 \\
2011 & 完全二叉树 & 5 & 单选 & 2 \\
2012 & 二叉树遍历 & 3 & 单选 & 2 \\
2013 & 二叉树遍历序列 & 5 & 单选 & 2 \\
2014 & 线索二叉树 & 5 & 单选 & 2 \\
2015 & 二叉树先后序序列 & 4 & 单选 & 2 \\
2017 & 树与二叉树的转换 & 6 & 单选 & 2 \\
2018 & 哈夫曼树 WPL & 5 & 单选 & 2 \\
2019 & 二叉树遍历 & 5 & 单选 & 2 \\
2020 & 哈夫曼编码 & 6 & 单选 & 2 \\
2021 & 二叉排序树 & 6 & 单选 & 2 \\
2022 & 线索二叉树前驱后继 & 6 & 单选 & 2 \\
2023 & 二叉树的遍历序列构造 & 5 & 单选 & 2 \\
2024 & 二叉树结点关系 & 6 & 单选 & 2 \\
\hline
\end{tabular}
\end{tablebox}

\noindent\textbf{高频：}二叉树的遍历序列构造（中序+先/后序→唯一二叉树），线索二叉树的前驱后继查找，哈夫曼树的 WPL 计算。

\subsection*{第5章 图}

\begin{tablebox}{图 — 真题索引（选择题）}
\centering
\begin{tabular}{|c|l|c|c|c|}
\hline
年份 & 知识点 & 题号 & 题型 & 分值 \\
\hline
2009 & 关键路径 & 5 & 单选 & 2 \\
2010 & 拓扑排序 & 6 & 单选 & 2 \\
2011 & 图的基本性质（连通图）& 7 & 单选 & 2 \\
2012 & 邻接矩阵存储性质 & 7 & 单选 & 2 \\
2013 & AOE 网关键路径 & 7 & 单选 & 2 \\
2014 & 图的遍历 & 7 & 单选 & 2 \\
2015 & 有向图顶点的度 & 5 & 单选 & 2 \\
2016 & 图的存储（邻接矩阵 vs 邻接表）& 7 & 单选 & 2 \\
2017 & 拓扑排序序列 & 7 & 单选 & 2 \\
2018 & 图的生成树 & 7 & 单选 & 2 \\
2020 & Dijkstra 算法 & 7 & 单选 & 2 \\
2021 & 最小生成树（Prim）& 7 & 单选 & 2 \\
2023 & 图的关键路径 & 7 & 单选 & 2 \\
2024 & 图的连通性 & 7 & 单选 & 2 \\
\hline
\end{tabular}
\end{tablebox}

\noindent\textbf{高频：}关键路径的 $ve$/$vl$ 计算、拓扑排序序列判定（DAG 才有）、图的存储方式的优缺点判断。

\subsection*{第6章 查找}

\begin{tablebox}{查找 — 真题索引}
\centering
\begin{tabular}{|c|l|c|c|c|}
\hline
年份 & 知识点 & 题号 & 题型 & 分值 \\
\hline
2009 & 平衡二叉树判断 & 4 & 单选 & 2 \\
2010 & 折半查找判定树 & 9 & 单选 & 2 \\
2011 & B-树（m 阶性质）& 10 & 单选 & 2 \\
2012 & B-树插入 & 9 & 单选 & 2 \\
2013 & 散列表 ASL（线性探测）& 9 & 单选 & 2 \\
2014 & 散列表 ASL & 9 & 单选 & 2 \\
2015 & B+树 & 9 & 单选 & 2 \\
2016 & B-树删除 & 10 & 单选 & 2 \\
2017 & B-树 & 10 & 单选 & 2 \\
2018 & B-树性质 & 10 & 单选 & 2 \\
2019 & B-树 & 10 & 单选 & 2 \\
2020 & 散列表 + 平均查找长度 & 10 & 单选 & 2 \\
2021 & B-树结点结构 & 9 & 单选 & 2 \\
2022 & B-树分裂与合并 & 9 & 单选 & 2 \\
2023 & B-树与B+树区别 & 9 & 单选 & 2 \\
2024 & 散列查找 & 9 & 单选 & 2 \\
\hline
\end{tabular}
\end{tablebox}

\noindent\textbf{规律：}B-树和散列表 ASL 是 408 查找部分的绝对核心。B-树的「$\lceil m/2\rceil$ 棵子树」等性质每年必考；散列表的探测次数计算也是高频。

\subsection*{第7章 排序}

\begin{tablebox}{排序 — 真题索引}
\centering
\begin{tabular}{|c|l|c|c|c|}
\hline
年份 & 知识点 & 题号 & 题型 & 分值 \\
\hline
2009 & 堆排序 & 7 & 单选 & 2 \\
2010 & 快速排序递归次数 & 7 & 单选 & 2 \\
2011 & 各排序算法比较 & 11 & 单选 & 2 \\
2012 & 快速排序 & 10 & 单选 & 2 \\
2013 & 堆排序建堆 & 11 & 单选 & 2 \\
2014 & 各排序趟的结果 & 11 & 单选 & 2 \\
2015 & 归并排序 & 11 & 单选 & 2 \\
2016 & 堆排序 & 11 & 单选 & 2 \\
2017 & 希尔排序 & 11 & 单选 & 2 \\
2018 & 排序算法选择 & 11 & 单选 & 2 \\
2019 & 堆排序（大根堆判定）& 11 & 单选 & 2 \\
2020 & 外部排序 & 11 & 单选 & 2 \\
2021 & 基数排序 & 11 & 单选 & 2 \\
2022 & 堆排序 & 11 & 单选 & 2 \\
2023 & 各排序算法稳定性 & 11 & 单选 & 2 \\
2024 & 堆排序 & 11 & 单选 & 2 \\
\hline
\end{tabular}
\end{tablebox}

\noindent\textbf{高频：}堆排序（建堆过程、大根堆判定）、各算法的稳定性判断、快速排序的趟数和递归深度。

\section{计算机组成原理}

\subsection*{第2章 数据的表示与运算}

\begin{tablebox}{数据表示 — 真题索引}
\centering
\begin{tabular}{|c|l|c|c|c|}
\hline
年份 & 知识点 & 题号 & 题型 & 分值 \\
\hline
2009 & 浮点数运算步骤 & 12,13 & 单选 & 2+2 \\
2010 & IEEE 754 表示（int/float 转换）& 18 & 单选 & 2 \\
2011 & IEEE 754 单精度 & 14 & 单选 & 2 \\
2012 & 补码范围、IEEE 754 & 13,14 & 单选 & 2+2 \\
2013 & 浮点数 IEEE 754 格式 & 13 & 单选 & 2 \\
2014 & IEEE 754 浮点数表示 & 13 & 单选 & 2 \\
2015 & IEEE 754 规格化 & 14 & 单选 & 2 \\
2016 & IEEE 754 特殊值 & 14 & 单选 & 2 \\
2017 & 补码与 IEEE 754 综合 & 13,14 & 单选 & 2+2 \\
2018 & IEEE 754 浮点数 & 14 & 单选 & 2 \\
2019 & IEEE 754 单精度 & 14 & 单选 & 2 \\
2020 & 补码与 IEEE 754 & 14 & 单选 & 2 \\
2021 & IEEE 754 浮点数 & 14 & 单选 & 2 \\
2022 & IEEE 754 & 14 & 单选 & 2 \\
2023 & 补码/IEEE 754 综合 & 13,14 & 单选 & 2+2 \\
2024 & IEEE 754 & 14 & 单选 & 2 \\
2025 & IEEE 754+补码运算 & 13,14 & 单选 & 2+2 \\
\hline
\end{tabular}
\end{tablebox}

\noindent\textbf{规律：}IEEE 754 浮点数表示每年 1-2 道选择题，是组成原理最稳定的考点。给定真值写编码/给定编码还原真值是核心能力。

\subsection*{第3章 存储系统}

\begin{tablebox}{存储系统 — 真题索引（选择题）}
\centering
\begin{tabular}{|c|l|c|c|c|}
\hline
年份 & 知识点 & 题号 & 题型 & 分值 \\
\hline
2009 & Cache 映射 & 15 & 单选 & 2 \\
2010 & Cache 组相联 & 15 & 单选 & 2 \\
2011 & Cache 直接映射地址划分 & 11 & 单选 & 2 \\
2012 & Cache 映射+替换 & 17 & 单选 & 2 \\
2013 & TLB+Cache+页表综合 & 16 & 单选 & 2 \\
2014 & Cache 映射 & 16 & 单选 & 2 \\
2015 & 虚拟存储器 & 17 & 单选 & 2 \\
2016 & Cache 替换算法（LRU）& 17 & 单选 & 2 \\
2017 & Cache+TLB 综合 & 17 & 单选 & 2 \\
2018 & Cache 映射方式 & 16 & 单选 & 2 \\
2019 & Cache 地址结构 & 17 & 单选 & 2 \\
2020 & Cache 映射（直接/组/全相联）& 17 & 单选 & 2 \\
2021 & 虚存+Cache 综合 & 17 & 单选 & 2 \\
2023 & 低位交叉编址 & 15 & 单选 & 2 \\
2024 & Cache+TLB 地址转换流程 & 15 & 单选 & 2 \\
\hline
\end{tabular}
\end{tablebox}

\noindent\textbf{高频综合题：}存储系统的大题几乎都是 Cache+页表+TLB 的综合题（2013, 2017, 2021），需要理解从虚拟地址→物理地址→Cache访问的完整流程。

\subsection*{第5章 CPU}

\begin{tablebox}{CPU — 真题索引}
\centering
\begin{tabular}{|c|l|c|c|c|}
\hline
年份 & 知识点 & 题号 & 题型 & 分值 \\
\hline
2009 & 流水线冒险 & 16 & 单选 & 2 \\
2010 & 流水线性能 & 19 & 单选 & 2 \\
2011 & 流水线数据冒险 & 20 & 单选 & 2 \\
2012 & 数据通路+指令执行 & 43,44 & 综合 & 13+13 \\
2013 & 流水线阻塞 & 18 & 单选 & 2 \\
2014 & 流水线转发 & 18 & 单选 & 2 \\
2015 & 数据通路分析 & 44 & 综合 & 13 \\
2016 & 流水线冒险 & 19 & 单选 & 2 \\
2017 & 五级流水线+转发 & 19 & 单选 & 2 \\
2018 & 流水线 CPI & 20 & 单选 & 2 \\
2020 & 数据冒险（RAW/WAW/WAR）& 20 & 单选 & 2 \\
2021 & 流水线吞吐率 & 20 & 单选 & 2 \\
2022 & 数据通路+控制信号 & 43 & 综合 & 13 \\
2023 & 流水线冒险类型 & 20 & 单选 & 2 \\
2024 & 流水线阻塞+转发 & 19 & 单选 & 2 \\
\hline
\end{tabular}
\end{tablebox}

\noindent\textbf{高频重点：}数据通路的控制信号分析（应用题）、三种流水线冒险的区分（RAW/结构/控制）、Load-use 冒险必须阻塞 1 拍。

\section{操作系统}

\subsection*{第2章 进程管理 — PV操作真题}

\begin{tablebox}{PV 操作与同步 — 真题索引}
\centering
\begin{tabular}{|c|l|c|c|c|}
\hline
年份 & 知识点 & 题号 & 题型 & 分值 \\
\hline
2009 & 生产者-消费者变体 & 45 & 综合 & 7 \\
2011 & 银行家算法（安全序列）& 46 & 综合 & 7 \\
2013 & PV 操作（读者-写者变体）& 46 & 综合 & 7 \\
2014 & PV 操作（理发师问题）& 46 & 综合 & 7 \\
2015 & PV 操作（独木桥问题）& 46 & 综合 & 7 \\
2017 & PV 操作（多缓冲区）& 46 & 综合 & 7 \\
2019 & PV 操作（抽烟者问题变体）& 45 & 综合 & 7 \\
2020 & 银行家算法 & 45 & 综合 & 7 \\
2023 & PV 操作综合 & 46 & 综合 & 7 \\
2024 & 信号量同步 & 45 & 综合 & 7 \\
\hline
\end{tabular}
\end{tablebox}

\subsection*{第3章 内存管理 — 真题}

\begin{tablebox}{内存管理 — 真题索引}
\centering
\begin{tabular}{|c|l|c|c|c|}
\hline
年份 & 知识点 & 题号 & 题型 & 分值 \\
\hline
2009 & 请求分页地址变换 & 46 & 综合 & 7 \\
2010 & 页式存储+TLB & 46 & 综合 & 7 \\
2011 & 页面置换（LRU/CLOCK）& 29,30 & 单选 & 2+2 \\
2012 & 请求分页+页面置换 & 46 & 综合 & 7 \\
2013 & 页表+TLB 综合 & 46 & 综合 & 7 \\
2014 & LRU 页面置换 & 30 & 单选 & 2 \\
2016 & 页面置换 CLOCK & 28,30 & 单选 & 2+2 \\
2017 & 页表地址变换 & 46 & 综合 & 7 \\
2019 & 页表综合 & 46 & 综合 & 7 \\
2020 & 地址变换+缺页 & 46 & 综合 & 7 \\
2022 & 页面置换+缺页率 & 46 & 综合 & 7 \\
2024 & TLB+页表+Cache+缺页综合 & 46 & 综合 & 7 \\
\hline
\end{tabular}
\end{tablebox}

\noindent\textbf{408 走向：}近年大题趋势是将\textbf{页表+TLB+Cache+缺页}串联成一条完整的访存流水线来考察，这对应了组成原理和操作系统的交叉考点。

\section{计算机网络}

\subsection*{第4章 网络层 — 真题}

\begin{tablebox}{网络层 — 真题索引}
\centering
\begin{tabular}{|c|l|c|c|c|}
\hline
年份 & 知识点 & 题号 & 题型 & 分值 \\
\hline
2009 & IP 分片 & 38 & 单选 & 2 \\
2010 & 子网划分 & 37,38 & 单选 & 2+2 \\
2011 & CIDR 路由聚合 & 37,38 & 单选 & 2+2 \\
2012 & IP 分片（片偏移计算）& 38 & 单选 & 2 \\
2013 & 路由表更新（RIP）& 35 & 单选 & 2 \\
2015 & ARP 协议 & 36 & 单选 & 2 \\
2016 & RIP 路由更新 & 37 & 单选 & 2 \\
2017 & IP 分片 & 38 & 单选 & 2 \\
2018 & 路由协议对比 & 37 & 单选 & 2 \\
2019 & CIDR 子网划分 & 37 & 单选 & 2 \\
2020 & IP 分片 & 38 & 单选 & 2 \\
2021 & CIDR 子网掩码 & 38 & 单选 & 2 \\
2023 & 路由聚合 & 37 & 单选 & 2 \\
2024 & ARP+IP 综合 & 47 & 综合 & 9 \\
\hline
\end{tabular}
\end{tablebox}

\noindent\textbf{高频：}IP 分片的片偏移计算、CIDR 的子网划分和路由聚合。网络层大题（47题）近年来多考 ARP+IP 的综合场景。

\subsection*{第5章 传输层 — 真题}

\begin{tablebox}{传输层 — 真题索引}
\centering
\begin{tabular}{|c|l|c|c|c|}
\hline
年份 & 知识点 & 题号 & 题型 & 分值 \\
\hline
2009 & TCP 拥塞控制（窗口变化图）& 39 & 单选 & 2 \\
2010 & TCP 序号与确认号 & 39 & 单选 & 2 \\
2011 & TCP 拥塞控制 & 39 & 单选 & 2 \\
2012 & TCP 连接管理 & 39 & 单选 & 2 \\
2013 & TCP 序号 & 38 & 单选 & 2 \\
2014 & TCP 拥塞控制（窗口图）& 38 & 单选 & 2 \\
2015 & TCP 三次握手 & 38 & 单选 & 2 \\
2016 & TCP 快速重传 & 39 & 单选 & 2 \\
2017 & TCP 拥塞避免 & 39 & 单选 & 2 \\
2018 & TCP 序号确认号 & 39 & 单选 & 2 \\
2019 & TCP 拥塞控制 & 39 & 单选 & 2 \\
2020 & TCP 三次握手+四次挥手 & 47 & 综合 & 9 \\
2021 & TCP 拥塞控制 & 39 & 单选 & 2 \\
2022 & TCP 序号+确认号 & 39 & 单选 & 2 \\
2023 & TCP 拥塞控制状态图 & 39 & 单选 & 2 \\
2024 & TCP 拥塞控制窗口图 & 47 & 综合 & 9 \\
\hline
\end{tabular}
\end{tablebox}

\noindent\textbf{规律：}TCP 拥塞控制几乎每年必考（选择题或综合题）。四种算法（慢开始/拥塞避免/快重传/快恢复）的拥塞窗口变化图要能默画。TCP 序号和确认号的计算也是高频。

\vspace{12pt}

\noindent\textbf{使用建议：}
\begin{enumerate}
    \item 每学完一章，对照本表找到对应年份真题，限时完成（选择题每题约 2 分钟，综合题每题约 15 分钟）。
    \item 真题答案和解析可参考 GitHub 仓库 \texttt{songmuhan/408}、\texttt{CodePanda66/CSPostgraduate-408} 或 csgraduates.com。
    \item 选择题优先完成近 10 年（2016-2025），综合题全部完成 2009-2025。
\end{enumerate}
